\documentclass{reportDoCS2025}

\usepackage{lmodern}
\usepackage{anyfontsize}

\addbibresource{biblio.bib}

%--------------------------------------
% \university{ĐẠI HỌC QUỐC GIA TP. HỒ CHÍ MINH}
% \school{TRƯỜNG ĐẠI HỌC KHOA HỌC TỰ NHIÊN}
\title{\MakeUppercase{Nghiệm trạng thái dừng của phương trình Naiver-Stokes}}
\subtitle{Lý thuyết Toán về Phương trình Navier-Stokes}

% \students{Lê \textsc{Nhựt-Nam} \\
%             Nhut \textsc{Nam Le} }

% \supervisors{Prof. Dr. Dumbledore \textsc{Albus} \\
% 	Prof. Dr. Snape \textsc{Severus} \\ 
%         Prof. Dr. McGonagall \textsc{Minerva} }

%--------------------------------------

\newcommand{\R}{\mathbb{R}}
\newcommand{\bv}{\mathbf{v}}
\newcommand{\bT}{\mathbf{T}}
\newcommand{\bD}{\mathbf{D}}
\newcommand{\bI}{\mathbf{I}}

\begin{document}

\cover
\numRoman

% \input{1.acknowledgements}

\contents
\listTables
\listImages


\pageNumber

\chapter*{\MakeUppercase{Mở đầu}}

Xem xét một chất lỏng nhớt (viscous fluid) có mật độ không thay đổi, hay còn gọi là \textit{viscous liquid} $\mathcal{L}$ di chuyển trong một miền cố định $\Omega \subseteq \R^3$.

Giả sử rằng một chuyển động tổng quát của $\mathcal{L}$ ứng với một hệ quy chiếu quán tính (inertial frame of reference) được mô tả bằng hệ phương trình sau:
\begin{equation}\label{eq:navier-stokes}
        \begin{aligned}
                \rho\left(\frac{\partial \bv}{\partial t} + \bv\cdot \nabla v\right) &= \mu\Delta \bv - \nabla \pi - \rho \mathbf{f}(x, t),\\
                \nabla \cdot v &= 0,
        \end{aligned}
\end{equation}
trong đó $t$ là thời gian, $x = (x_1, x_2, x_3)$ là một điểm của $\Omega$, $\rho$ là mật độ của $\mathcal{L}$ (hằng số), $\bv = \bv(x, t) = (v_1(x, t), v_2(x,t), v_3(x, t))$ và $\pi = \pi(x, t)$ tương ứng là trường vận tốc và trường áp suất. Và hằng số dương $\mu$ đại diện cho hệ số nhớt (viscosity coefficient).

Hơn nữa, thành phần đối lưu (convective term) là:
\begin{equation}
        \bv \cdot \nabla \bv \equiv \sum_{i=1}^3 v_i\frac{\partial \bv}{\partial x_i},
\end{equation}

Trong các phần bên dưới, ta sử dụng một số toán tử. Toán tử Laplace được định nghĩa như sau:
\begin{equation}
        \Delta \equiv \sum_{i=1}^3\frac{\partial^2}{\partial x_i^2},
\end{equation}
và toán tử gradient được định nghĩa như sau:
\begin{equation}
        \nabla \equiv \left(\frac{\partial}{\partial x_1}, \frac{\partial}{\partial x_2}, \frac{\partial}{\partial x_3}\right).
\end{equation}

Còn phân kỳ của $\bv$ (hay divergence), được định nghĩa:
\begin{equation}
        \nabla \cdot \bv \equiv = \sum_{i = 1}^3\frac{\partial v_i}{\partial x_i}
\end{equation}

Cuối cùng, $-\mathbf{f}$ là ngoại lực trên mỗi đơn vị khối lượng (body force) tác động lên $\mathcal{L}$.

Trong ngôn ngữ của cơ học thuần túy hiện đại (rational mechanics), dưới giả định cấu thành đơn giản trên chất lưu $\mathcal{L}$ mà hình thành nên hệ~\eqref{eq:navier-stokes} là một thành phần động của \textit{tensor ứng suất (Cauchy stress tensor)} $\bT$ và có mối liên hệ tuyến tính với \textit{tensor kéo dãn (stretching tensor)} $\bD$:
\begin{equation}\label{eq:cauchy-stress}
        \bT = -\pi\bI + 2\mu\bD,
\end{equation}
trong đó, $\bI$ là ma trận đơn vị và $\bD = \{D_{ij}\}$ với
\begin{equation}
        D_{ij} = \frac{1}{2}\left(\frac{\partial v_i}{\partial x_j} + \frac{\partial v_j}{\partial x_i}\right).
\end{equation}
Chất lỏng mà thỏa phương trình~\eqref{eq:cauchy-stress} cũng được gọi là \textit{Newtonian liquid hay chất lưu Newton}.

Trong bài toán liên quan đến tính duy nhất nghiệm của phương trình\eqref{eq:navier-stokes}, nghiệm hai chiều mô tả \textit{chuyển động mặt phẳng} của $\mathcal{L}$ nhận được nhiều sự chú ý. Với loại nghiệm này, trường $\bv$ và $p$ chỉ phụ thuộc vào $x_1, x_2$, và $t$ và hơn nữa $v_3 \equiv -$. Hệ quả là miền liên quan của dòng chảy $\Omega$ trở thành một tập con của mặt phẳng $\R^2$. 

Mục tiêu chính trong nghiên cứu phương trình\eqref{eq:navier-stokes} là tìm hiểu những tính chất cơ bản và có ý nghĩa của $\bv$ và $\pi$ như sự tồn tại nghiệm, tính duy nhất nghiệm, tính chính quy, và đặc tính tiệm cận trong không gian (khi mà $\Omega$ không bị chặn) và thời gian. Các tính chất này có thể khác nhau khi chúng ta xem xét $\mathcal{L}$ trong trạng thái dừng (steady state) hay không dừng (unsteady state), trường vận tốc và áp suất phụ thuộc hay không phụ thuộc vào thời gian.   

Trong phần tóm tắt này, chúng ta xem xét chuyển động dừng đối với hệ\eqref{eq:navier-stokes} với $\partial \bv \setminus \partial t \equiv -$ và $\mathbf{f} = \mathbf{f}(x)$ có dạng như sau:
\begin{equation}\label{eq:steady-navier-stokes}
        \begin{aligned}
                \nu\Delta\bv &= \bv \cdot \nabla \bv + \nabla p + \mathbf{f},\\
                \nabla \cdot \bv &= 0,
        \end{aligned}
\end{equation}
trong đó, $p = \pi \setminus \rho$ và $\nu = \mu \setminus \rho$ là hệ số độ nhớt động học (kinematic viscosity coefficient). Ta có thể gọi $p$ là ``áp suất'' (hay trường áp suất) lên chất lưu $\mathcal{L}$.



%--------------------------------------
\chapter{\MakeUppercase{Dòng chảy trong miền bị chặn}}





%--------------------------------------
\chapter{\MakeUppercase{Dòng chảy trong miền bên ngoài}}


\section{Dòng chảy ba chiều}



\section{Dòng chảy mặt phẳng}



%--------------------------------------
\chapter{\MakeUppercase{Dòng chảy trong miền có biên không bị chặn}}


% \textbf{Cite a source} from \texttt{biblio.bib} with \texttt{\textbackslash cite}: \cite{example_of_source}.\\

% \textbf{Display a bibliography} quickly with \texttt{\textbackslash insertbibliography}.
% \insertbibliography

\end{document}